\section*{Simulation Results and Discussion}\label{ch:ch5label}

Let us first try only with cash-in-advance, money flow and normalizing (c and l being between 0 and 1) constraints. In this case, we are just trying to maximize for a single person. Let us also assume that the individual begins with a fair amount of money (0.5 units) so that he does not starve in the first time step.

The result: c = [1,1,...,1] and l = [0,0,...,0] regardless of the interest rate, prices, discount factor and any other variable. It is a bit weird that the consumer is able to consume a lot without actually working. What we notice is that amount of bond bought is [3.5,3,...] yet the debts are not paid. So a Ponzi scheme is created. The consumer borrows more and more each round. Even though not being an economically valid result, let us analyze in terms of convex optimization terminology. Currently, the only active constraints are normalizing constraints (i.e. $0 \leq c \leq 1$). So the Lagrange multipliers for those are strictly positive and the rest are zero.

Now let us try adding ponzi scheme preventing constraint in the last time step, s (amount of bonds) should be zero since the consumer will not be existent in next time step.

$s_{N-1} = 0$ 

In this case, R, $\beta$ and other modeling parameters begin to matter. Let us include a 3D graph of c(t) versus $\beta$.

\begin{figure}
\caption{Optimal consumption strategy for various $\beta$ and $m(1)=0.5$, R=0, no market clearing, Ponzi not allowed}
\centering
\includegraphics[width=12cm,height=8cm]{figures/r0_consumption.pdf}
\end{figure}


\begin{figure}
\caption{Optimal labor strategy for various $\beta$ and $m(1)=0.5$, R=0, no market clearing, Ponzi not allowed}
\centering
\includegraphics[width=12cm,height=8cm]{figures/r0_labor.pdf}
\end{figure}


\begin{figure}
\caption{Optimal consumption strategy for various $\beta$ and $m(1)=0.5$, R=0.8, no market clearing, Ponzi not allowed}
\centering
\includegraphics[width=12cm,height=8cm]{figures/r08_consumption.pdf}
\end{figure}


\begin{figure}
\caption{Optimal labor strategy for various $\beta$ and $m(1)=0.5$, R=0.8, no market clearing, Ponzi not allowed}\centering
\includegraphics[width=12cm,height=8cm]{figures/r08_labor.pdf}
\end{figure}


As it can be seen (on graphs following this page), for low $\beta$, the solution found is to maximize initial consumption and to minimize initial labor. The mechanism which allows this behavior is interest rate (R). As beta rises, we see that consumption begins to spread more evenly. 

Now let us analyze the dual variables and the values they attain. From those variables, we can comment on duality of the problem and also on the which constraints are active.

Again for the same R (=0.8) and same initial income (m=0.5) let us now plot the duality variables and afterwards let us comment on them.

\begin{figure}
\caption{Dual variable $\lambda$ for varying $\beta$ and $m(1)=0.5$, R=0.8, no market clearing, Ponzi not allowed}
\centering
\includegraphics[width=12cm,height=8cm]{figures/r08_lambda.pdf}
\end{figure}


\begin{figure}
\caption{Dual variable $\sigma$ for various $\beta$ and $m(1)=0.5$, R=0, no market clearing, Ponzi not allowed (errata: X axis is $\beta$ and not time!)}
\centering
\includegraphics[width=12cm,height=8cm]{figures/r08_sigma.pdf}
\end{figure}


\begin{figure}
\caption{Dual variable $\rho$ for various $\beta$ and $m(1)=0.5$, R=0.8, no market clearing, Ponzi not allowed}
\centering
\includegraphics[width=12cm,height=8cm]{figures/r08_rho.pdf}
\end{figure}


\begin{figure}
\caption{Dual variable $\epsilon$ for various $\beta$ and $m(1)=0.5$, R=0.8, no market clearing, Ponzi not allowed}\centering
\includegraphics[width=12cm,height=8cm]{figures/r08_epsilon.pdf}
\end{figure}

Important results can be obtained from these graphs of dual variables. Let us first denote which dual variable corresponds to which constraint and at the same time discuss the graphs obtained:

\begin{itemize}
\item $\epsilon: 0\leq c \leq 1$ \\

for different $\beta$s, this variable is always 0. This means that for high interest rates, we are always inside this boundary and never at the boundary. Only for very low  beta values, we begin to come at $0\leq c$ boundary, since with bonds the consumer chooses to spend early and work afterwards to compensate. This dual variable symbolizes how "unworthy" is a consumption good. 

\item $\rho: 0\leq l \leq 1$ \\

for different $\beta$s, this variable is very high at initial times then gets low. This means that the incentive is to work less when it matters for the utility, afterwards it becomes zero. This dual variable symbolizes how "worthy" is leisure, (i.e. a shadow price) 

\item $\lambda: m(t+1) + s(t+1) = m(t) + (1+R)s(t) + p(k)c(k) + tau(k) - p(k)c(k)$\\

this variable decreases with time and increases with beta. It is not wrong to consider this one as an incentive to "print money". The reason is, lambda is the dual variable weighing the constraint of "earned money is equal to spent money". As this dual variable gets away from zero, we see that the consumer is in demand of more money but is not able to get since he did not work hard enough.

\item $\sigma: s(N-1)=0 $\\

this variable increases with beta. This constraint was on "to leave debt for the offsprings". It makes sense since as the beta increases, the consumption at later life stages begins to be important for the individual, hence the incentive to "keep the spending trend as is." So $\sigma$ is the incentive of "leaving debt to offsprings".

\end{itemize}